\chapter{ماهیت داده}
\label{ch:nature-of-data}

این فصل سوالی رو مطرح می‌کند که اکثر مردم از پاسخ به سوال صرف نظر می‌کنند و تظاهر می‌کنند از قبل می‌دانند آن چیست: داده چیست؟
معمولاً این سؤال با پاسخ‌هایی بسیار خلاصه رد می‌شود و بحث سریع به سراغ جاهای «کاربردی‌تر» می‌رود. مانند:
داده چگونه ذخیره می‌شود، چگونه جستجو می‌شود، چگونه فشرده می‌شود، و دیگر مسائل کاربردی دیگر. پاسخ بیش از آنچه در ابتدا به نظر می‌رسد اهمیت دارد. نحوه‌ی تفکر شما در مورد داده‌ها و هر تصمیمی که هنگام کار با داده‌ها می‌گیرید را شکل می‌دهد.
\newpage
\section{یک سؤال که همه فکر می‌کنند پاسخش را می‌دانند}
\label{sec:question-everyone-thinks-they-know}

اگر بخواهیم به تاریخ بشریت نگاهی بیندازیم می‌بینیم که در ۲۰-۳۰ سال گذشته، بر خلاف بیشتر طول تاریخ بشر، انسان به دسترسی تقریباً آنی داده رسیده، و تقریباً پاسخ به هر سوالی، تنها چند نانوثانیه فاصله دارد (با شروع عصر هوش مصنوعی می‌توانیم بگوییم این سرعت افزایش پیدا کرده).

یا به قول مارکس:

در گذشته سنت یا نسل‌های گذشته مانند کوهی بر مغز زندگان سنگینی می‌کرد \lr{(THE EIGHTEENTH BRUMAIRE --- OF --- LOUIS BONAPARTE.)}. اما در طول این ۲۰-۳۰ سال گذشته، این دسترسی بی‌سابقه، پارادوکسی را هم با خودش آورده. بسیاری از اعضای آنچه نسل هزاره و نسل Z نامیده می‌شود، شروع به زیرسوال بردن ارزش‌ها، هدف‌ها، عقاید و... کرده‌اند. چون اگر پاسخ همیشه در دسترس است، دیگر چه نیازی به رنج کشیدن برای رسیدن به آن هست؟

و خب در دسترس بودن همه‌ی جواب‌ها پیش از پرسیدن سوال، باعث ایجاد پارادوکسی عجیب و غیرمنتظره می‌شود، که مسئله این‌طور جا افتاده که «داده زیاد شده و بد هست»، اما مسئله این هست که \textbf{در دسترس بودن داده با در اختیار داشتن خرد، دو چیز کاملاً متفاوت‌اند.} به عنوان مثال من می‌توانم برای یک شخص رندوم، بدون اینکه بدونم تاریخ فرانسه چیست، هزاران سند معتبر پیدا کنم؛ بدون اینکه بدونم این تاریخ درباره‌ی چی حرف می‌زنه و چرا این تاریخ به این شکل رقم خورده. در ادامه می‌توانیم این گپ بزرگ بین فهمیدن و داشتن متن‌های خوش‌خط‌وخال را دسته‌بندی کنیم.
\newpage
\subsection{بازگشت به سؤال اصلی}
\label{sub:return-to-main-question}

پس بیایید به سؤال اصلی برگردیم: داده چیست؟ داده را می‌توان تقریباً هر چیزی در نظر گرفت. برای آنکه این ادعا انتزاعی نماند، بیایید آن را روی یک متن واقعی امتحان کنیم. به عنوان مثال در کتاب \textit{\lr{The Eighteenth Brumaire of Louis Bonaparte}} نوشته‌ی کارل مارکس:

\begin{latin}
\begin{quote}
\itshape
``Upon the overthrow of the first Napoleon came the restoration of the Bourbon throne (Louis XVIII, succeeded by Charles X). In July, 1830, an uprising of the upper tier of the bourgeoisie, or capitalist class --- the aristocracy of finance ---, overthrew the Bourbon throne, or landed aristocracy, and set up the throne of Orleans, a younger branch of the house of Bourbon, with Louis Philippe as king. From the month in which this revolution occurred, Louis Philippe's monarchy is called the ``July Monarchy.'' In February, 1848, a revolt of a lower tier of the capitalist class --- the industrial bourgeoisie ---, against the aristocracy of finance, in turn dethroned Louis Philippe. This affair, also named from the month in which it took place, is the ``February Revolution.'' The ``Eighteenth Brumaire'' starts with that event.

Despite the inapplicableness to our own affairs of the political names and political leadership herein described, both these names and leaderships are to such an extent the products of an economic-social development that has here too taken place with even greater sharpness, and they have their present or threatened counterparts here so completely, that, by the light of this work of Marx', we are best enabled to understand our own history, to know whence we came, and whither we are going, and how to conduct ourselves.''
\end{quote}
\end{latin}

پیش از آنکه ادامه‌ی متن را بخوانید، از خودتان بپرسید: شما دقیقاً چه اطلاعات و داده‌ای از این متن دریافت می‌کنید؟

اگر بخواهیم بسیار کلی‌نگر باشیم، می‌توانیم سرنوشت ناپلئون اول و دوم، و همچنین انقلاب‌های فرانسه را از این متن دریافت کنیم. اگر دقیق‌تر نگاه کنیم، می‌توانیم اطلاعاتی مثل نام پادشاهان، تاریخ‌های مشخص (۱۸۳۰، ۱۸۴۸)، و حتی طبقه‌بندی‌های اجتماعی (اشرافیت مالی، بورژوازی صنعتی) را از همین چند خط استخراج کنیم.

اما موضوع واقعاً جالب می‌شود وقتی زاویه‌ی دید خواننده را تغییر دهیم.

\subsubsection{یک کودک}
که این متن را بخواند، تنها داده‌های بسیار ساده‌تری تشخیص می‌دهد: «این متن انگلیسی است، درباره‌ی آدم‌هایی به نام ناپلئون و لوئی صحبت می‌کند، و چند عدد سال مثل ۱۸۳۰ و ۱۸۴۸ در آن هست.»

\subsubsection{یک خواننده‌ی آشنا با زبان انگلیسی}
ممکن است داده‌های زبانی بیشتری دریافت کند: «این متن به زبان انگلیسی نوشته شده، از ساختارهای نحوی خاصی استفاده می‌کند، و واژگان تاریخی و سیاسی دارد.»

\subsubsection{یک مورخ}
ممکن است چیز دیگری ببیند: «این متن درباره‌ی تحولات سیاسی فرانسه، در فاصله‌ی سقوط ناپلئون تا انقلاب ۱۸۴۸، است و جابه‌جایی قدرت میان گروه‌های اجتماعی و سیاسی مختلف را توصیف می‌کند.»

\subsubsection{یک اقتصاددان سیاسی}
به‌جای تمرکز بر نام پادشاهان، به ساختار طبقاتی و رابطه‌ی میان گروه‌های اقتصادی توجه می‌کند: نقش بورژوازی، اشرافیت مالی، و بورژوازی صنعتی.

\subsubsection{یک فیلسوف سیاسی}
متن را به‌عنوان داده‌ای درباره‌ی قدرت، طبقه، دولت، انقلاب، و تغییر ساختارهای اجتماعی می‌بیند.

\subsubsection{یک پژوهشگر ادبی}
به سبک نگارش، واژگان، ساختار روایت، و شیوه‌ی بازنمایی تاریخ در متن توجه می‌کند.

از همین متن، همچنین می‌توان تقسیم قدرت میان طبقه‌ی پرولتاریا و بورژوازی در فرانسه را دریافت کرد، و خب همان‌طور که گفته شد می‌شود فهمید این متن اثری از مارکس، فیلسوف و نظریه‌پرداز آلمانی، هست.

نکته‌ی بعدی که می‌توان از متن به دست آورد نکات تحلیلی هست که مستقیماً از متن نوشته نشده‌اند، بلکه از تحلیل متن به دست می‌آیند. برای مثال عبارت:

\begin{latin}
\begin{quote}
``Louis Philippe's monarchy is called the `July Monarchy'''
\end{quote}
\end{latin}

به‌صورت مستقیم یک رابطه میان \textit{Louis Philippe} و \textit{July Monarchy} را بیان می‌کند، رابطه‌ای که خود متن آن را صریحاً اعلام نکرده، بلکه خواننده آن را با تحلیل و بررسی متوجه شده.

حال با فهمیدن این نکات باید به سراغ نکته‌ی کانونی برویم. تمام این‌ها از یک رشته‌ی ثابت و کاراکترهای یکسان بیرون آمده‌اند. چیزی در خود متن تغییر نکرده، و هر خواننده بر اساس دانش خود متن را قضاوت می‌کند.

\section{هرم DIKW: از رشته‌ی خام تا خرد}
\label{sec:levels-of-data}

نکات و چیزی که در بالا دیدیم، در واقع یک سلسله‌مراتب از سطوح مختلف است. چیزی که در ادبیات علم داده و مدیریت اطلاعات با عنوان استاندارد هرم \textbf{DIKW} شناخته می‌شود، که به‌صورت داده، اطلاعات، دانش، و خرد است.

این هرم دقیقاً همان فاصله و گپی هست که در ابتدا اشاره شد، و پارادوکس عجیب در بین نسل‌های جدید بشریت هست که نشون می‌ده داشتن داده‌های زیاد هیچ تضمینی به رسیدن به لایه‌ی چهارم و خرد نمی‌دهد.
\newpage
\subsection{هرم DIKW چیست؟}
\label{sub:what-is-dikw}

\subsubsection{Data}
چیزی که مستقیماً در اختیار ماست: رشته‌ای از کاراکترها یا کلمات، بدون هیچ تفسیر یا سازمان‌دهی اضافه. در مثال بالا، خودِ متن انگلیسیِ نقل‌شده  پیش از هرگونه تحلیل داده‌ی خام است. این همان لایه ای هست که اشاره شد و در کسری از ثانیه در دسترس هرکسی هست.

\subsubsection{Information}
زمانی که داده‌های خام از متن استخراج و در قالبی مشخص سازمان‌دهی می‌شوند، و در یک ساختار و زمینه‌ی معین معنا پیدا می‌کنند. برای مثال، «انقلاب فوریه‌ی ۱۸۴۸» یا «سلطنت ژوئیه» اطلاعات هستند: داده‌ی خام سازمان‌یافته در قالب رویدادهای تاریخی مشخص.

\subsubsection{Knowledge}
زمانی که ترکیب اطلاعات و روابط میان آن‌ها، درکی درباره‌ی یک موضوع به ما می‌دهد. درک تحول سیاسی و طبقاتی فرانسه در این دوره، و سرنوشت فرانسه پس از سقوط ناپلئون، دانش است. نه صرفاً فهرستی از تاریخ‌ها و نام‌ها، بلکه فهمی از \emph{چرایی} و \emph{رابطه‌ی} میان آن‌ها.

\subsubsection{Wisdom}
بالاترین لایه‌ی این هرم، جایی است که دانش به قضاوت و عمل درست تبدیل می‌شود: نه صرفاً دانستنِ اینکه چرا سلطنت ژوئیه سقوط کرد، بلکه توانِ به‌کار بردنِ آن فهم برای خواندنِ درستِ رویدادهای مشابه، امروز. همان چیزی که در ابتدای فصل غایب بود: دسترسیِ آنی به هزاران سند تاریخی (داده)، خودش هرگز جای خردِ لازم برای فهمیدنِ \emph{چرا} یک رویداد اهمیت دارد را نمی‌گیرد. درواقع میتوانیم بگویم تفاوت میان افراد، سازمان‌ها و سیستم‌ها اغلب نه در مقدار داده‌ای که در اختیار دارند، بلکه در توانایی آن‌ها برای سازمان‌دهی، تفسیر و استفاده از آن داده‌ها آشکار می‌شود.

\section{آنچه هیچ تعریفی مستقیماً نمی‌گوید}
\label{sec:what-no-definition-says}

اکنون که جایگاه داده و مسیر رسیدن به خرد را درک کردیم، بیایید پیش از هر نتیجه‌گیری و پیش‌قضاوت، نگاهی به ریشه‌ی خود این کلمه یعنی data بیندازیم.

به‌صورت تاریخی، کلمه‌ی data از کلمه‌ی لاتین datum، اسم مفعول فعل dare به معنای دادن، گرفته شده هست.

ما در گفتار روزمره شکل data را اغلب به‌عنوان اسم جمع مفرد در نظر می‌گیریم («داده‌ها واضح هستند»)، اما در نوشتار فنی صرف جمع ترجیح داده می‌شود.

از آنجا که هر فردی می‌تواند برای هر متنی قضاوت خودش را داشته باشد، خود کلمه‌ی data هم استثنا نیست، و در سطح جهانی هر فرد یا سازمانی، داده را از منظر خودش تعریف کرده:

\textbf{راب، موریس و کورونل} آن را «حقایق خام» می‌دانند که هنوز پردازش نشده‌اند تا معنایشان آشکار شود؛
\textbf{فرهنگ لغت وبستر} آن را چیزی «داده‌شده یا پذیرفته‌شده» می‌داند، مبنایی که یک استنتاج از آن حاصل می‌شود؛
\textbf{فرهنگ لغت آکسفورد}، مختصرتر، آن را حقایق یا چیزهای شناخته‌شده‌ای می‌داند که مبنای استنتاج یا محاسبه قرار می‌گیرند.

به‌صورت کلی و روی‌هم‌رفته، تمام این تعریف‌ها به چیزی اشاره دارند که هیچ‌کدام آن را مستقیماً بیان نمی‌کنند:

\begin{quote}
\itshape
داده‌ها هیچ شکلی ندارند که مستقل از یک دیدگاه خاص به آن‌ها تعلق داشته باشد. یک واقعیت فقط زمانی به‌عنوان داده محسوب می‌شود که در یک چارچوب، چارچوبی از معنا، هدف، یا کاربرد، قرار گیرد که نسبت به آن مرتبط، سازمان‌یافته، منسجم، و مفید باشد.
\end{quote}
\newpage
یک حلقه‌ی رشد، برای سیستمی که مفهومی از فصول ندارد، هیچ معنایی ندارد؛ یک ولتاژ، برای مداری که برای نادیده‌گرفتنِ آن ساخته شده، هیچ معنایی ندارد. داده‌ها نوعی توافقِ قبلی، هرچند غیررسمی، درباره‌ی آنچه که در وهله‌ی اول ارزشِ توجه دارد را پیش‌فرض می‌گیرند. این همان نکته‌ای است که مثال مارکس آن را به‌روشنی نشان داد: کودک، مورخ، و اقتصاددان، هر سه به یک رشته‌ی یکسان از کاراکترها نگاه کردند، اما سه «داده»ی متفاوت دریافت کردند، چون سه چارچوب متفاوت با خود آوردند.

اما یک چیزی در همان متن مارکس هنوز بی‌جواب مونده هست. سلطنتی که هیچ ارتباط ذاتی به ماه ژوئیه نداشت، سلطنت ژوئیه نامیده شد؛ فقط به این دلیل که در آن ماه آغاز شد. آیا دلیل خوبی بود؟ مانند انقلابی که ربطی به مفهوم فوریه نداشت، انقلاب فوریه خوانده شد. چرا نام دیگری برای آن انتخاب نشد؟ یک اسم برای چیزی که خودش نیست به کار رفت، و همه بی هیچ اعتراضی همان اسم را پذیرفتند و آن انقلاب را به آن اسم به رسمیت شناختند.

درواقع ما می‌توانیم بگوییم این اسم‌ها چیزی را در خود متن حمل نمی‌کردند؛ آن‌ها تنها به‌عنوان نشانه‌ای برای اشاره به چیزی دیگر، بیرون از خودشان، به کار گرفته شده بودند. حال آن سوال بی‌جواب چیست؟

چگونه یک کلمه یا عدد یا نشانه می‌تونه این‌گونه جای چیزی دیگر را بگیرد؟ 

این پرسش، ما را از خودِ داده به مسئله‌ای بنیادی‌تر می‌رساند: رابطه میان چیزی که می‌بینیم و چیزی که آن را نمایندگی می‌کند.